\documentclass[10pt,a4paper]{amsart}
\usepackage[margin=19mm]{geometry}
\usepackage{amsmath,amssymb,mathtools}
\usepackage[scr=rsfs,cal=euler]{mathalfa}
\usepackage{xfrac}
\usepackage[unicode,hidelinks,bookmarks=false]{hyperref}
\newcommand{\ie}{i.e.\ }
\newcommand{\cf}{cf.\ }
\newcommand{\resp}{respectively\ }
\newcommand{\Subsection}{\subsection}
\providecommand{\nobreakdash}{\nobreak\hbox{-}}
\setlength{\emergencystretch}{2em}
\multlinegap0pt
\usepackage{amssymb}
\usepackage{mathtools}
\usepackage{float}
\usepackage{tikz}
\usepackage[shortlabels]{enumitem}
\newtheorem{theorem}{Theorem}
\newtheorem{corollary}{Corollary}
\newcommand{\ceil}[1]{\left\lceil #1 \right\rceil}
\title{Extremal chromatic bounds for distance Laplacian eigenvalues}
\author{Bilal Ahmad Rather}
\date{Source: arXiv:2604.10785v2 (CC BY 4.0)}
\begin{document}
\maketitle

\noindent\emph{Source and licence.} Extremal chromatic bounds for distance Laplacian eigenvalues. Version arXiv:2604.10785v2 (CC BY 4.0), distributed by arXiv under the Creative Commons Attribution 4.0 International licence, http://creativecommons.org/licenses/by/4.0/, as recorded in the arXiv licence metadata for this exact version and shown on the abstract page https://arxiv.org/abs/2604.10785v2. Full licence notice: \texttt{licenses/arxiv-2604.10785-CC-BY-4.0.txt}.\par\smallskip

\noindent\textit{Editorial scope.} Counted unit: the article's corollary cor:many-above-bchi (source0.tex lines 334--346). The article's complete original proof of it is delivered with it (source0.tex lines 348--401, one \texttt{proof} environment of 54 lines). No mathematics is written, repaired or re-derived: the statement and the proof are the article's own lines, in the article's own notation, and no retained line is substituted or re-flowed.  Retained uncounted as the source-local context the proof needs: lines 258--264.  Nothing was omitted from any retained span, so the package carries no omission record.  No retained line contains a citation, so the wrapper carries no bibliography and no bibliography entry.  No retained line contains a figure environment, an image-inclusion command or drawing code, so no image is delivered and the package declares no asset dependency.  Editorial formatting only, in addition: an AMS \texttt{amsart} preamble that restates the article's own declarations and macros used by the retained lines, namely the environments \texttt{theorem}, \texttt{corollary} on separate counters, together with the standard packages those lines need.  The retained environments print in this document as Theorem 1, Corollary 1 in order of appearance, whereas the article prints the same items with the numbering of its own full document.  The complete native source of the article as distributed in the arXiv e-print for this exact version is bundled unchanged as \texttt{sources/198293/source0.tex}.
\begin{theorem}[Color-class majorization]\label{thm:color-majorization}
Let $G$ be connected, let $\chi=\chi(G)$, and let
$\ell_1\ge\cdots\ge\ell_\chi$ be the color-class sizes of a fixed optimal $\chi$-coloring. Then
 $\partial^{L}_i(G)\ge n+\ell_1$ for $1\le i\le \ell_1-1. $
More generally, let $s_j=\sum_{t=1}^{j}(\ell_t-1)$ for $j\ge1,$ and $s_0=0,$ and let $p=|\{j:\ell_j\ge2\}|$. Then
$\partial^{L}_i(G)\ge n+\ell_j$  whenever $s_{j-1}+1\le i\le s_j,$ for  $1\le j\le p.$ 
\end{theorem}

\begin{corollary}\label{cor:many-above-bchi}
For every connected graph $G$ with $\chi(G)=\chi\le n-1$,
\[
    \left|\left\{i\in\{1,
\dots,n\}:\partial_i^L(G)\ge b_\chi\right\}\right|
    \ge \ell_1-1\ge \ceil{\frac{n}{\chi}}-1.
\]
Equivalently,
\[
    \mu_G([0,b_\chi))\le n-\left(\ceil{\frac{n}{\chi}}-1\right)
    =n-\ceil{\frac{n}{\chi}}+1.
\]
\end{corollary}

\begin{proof}
	Let $V(G)=C_1\dot\cup C_2\dot\cup\cdots\dot\cup C_\chi$ be the fixed optimal $\chi$-coloring, with
	\[
	\ell_1=|C_1|\ge \ell_2=|C_2|\ge\cdots\ge \ell_\chi=|C_\chi|.
	\]
	By Theorem~\ref{thm:color-majorization}, the first $\ell_1-1$ distance Laplacian eigenvalues satisfy
	 $\partial_i^L(G)\ge n+\ell_1$ for  $1\le i\le \ell_1-1.$ Since the color classes form a partition of the $n$ vertices into $\chi$ nonempty sets, their average size is $n/\chi$. Therefore the largest color class has size at least the ceiling of this average, that is,
	\begin{equation}\label{eq:ell1-pigeonhole}
		\ell_1\ge \ceil{\frac{n}{\chi}}.
	\end{equation}
	It follows that, for every $1\le i\le \ell_1-1$, we have
	\[
	\partial_i^L(G)
	\ge n+\ell_1
	\ge n+\ceil{\frac{n}{\chi}}
	=
	b_\chi.
	\]
	Thus, the eigenvalues $\partial_1^L(G),\partial_2^L(G),\dots,\partial_{\ell_1-1}^L(G)$ 	all lie in the interval
	 $[b_\chi,\partial_1^L(G)].$ 	Hence, we have  $\left|\left\{i\in\{1,\dots,n\}:\partial_i^L(G)\ge b_\chi\right\}\right|
	\ge \ell_1-1.$ By \eqref{eq:ell1-pigeonhole}, we also have
	\[
	\ell_1-1\ge \ceil{\frac{n}{\chi}}-1.
	\]
	 It remains only to translate this into the equivalent upper bound for the number of eigenvalues below $b_\chi$. Since the eigenvalues are counted with algebraic multiplicity and are ordered from largest to smallest, every eigenvalue lies either in
	 $[0,b_\chi)$ or in $[b_\chi,\partial_1^L(G)].$ 
	Therefore, we obtain
	\[
	\mu_G([0,b_\chi))
	+
	\mu_G([b_\chi,\partial_1^L(G)])
	=
	n.
	\]
	The inequality just proved gives
	\[
	\mu_G([b_\chi,\partial_1^L(G)])
	\ge
	\ell_1-1
	\ge
	\ceil{\frac{n}{\chi}}-1.
	\]
	Consequently, we obtain
	\[
	\mu_G([0,b_\chi))
	=
	n-\mu_G([b_\chi,\partial_1^L(G)])
	\le
	n-\left(\ceil{\frac{n}{\chi}}-1\right).
	\]
	Thus, we have  $\mu_G([0,b_\chi))
	\le
	n-\ceil{\tfrac{n}{\chi}}+1$ 
\end{proof}

\end{document}
